\scanpage{173}
% archive sheet number: 172
est bijectif. Ceci résulte de qqe chose de
plus général, voir plus bas.

Ainsi, on a identification des \uline{groupes} $G$
avec une algèbre $A_0$ de fonctions de $G$ dans $k$,
satisfaisant aux conditions précédentes, et
\uline{certains} \deleted{$\unsure{\underline U}$} $K$-anneaux $\underline U$

\uline{N.B.} Soit $\underline U$ un anneau \inserted{\unsure{linéairemt|linéairement}} \unsure{topologisé|topologique}
complet ayant un syst. fond. de
voisinages de $0$ \inserted{$\unsure{I}$} tels que les $U/I$ sont
de dim finie, $A$ la coalgèbre duale
(réunion \unsure{filtrante|filtrée} de ses sous-coalgèbres de
dim finie). Les $\underline U$-modules
permis sont les modules unitaires
sur $A$ de dim. finie sur $k$.

Soit $\Omega$ une \deleted{\unsure{bialgèbre}} \inserted{\unsure{$k$-algèbre|algèbre}} associative avec
$1$, considérons la classe des $\Omega$-modules
unitaires \deleted{Soit} de dim finie sur $k$.
Soit $A$ l'ensemble des formes sur $\Omega$
coefficients de représentations linéaires
de \uline{dim. finie} de $\Omega$. La multiplication
de $\Omega$ est continue pour $\sigma(\Omega\otimes\Omega,A\otimes A)$
et $\sigma(\Omega,A)$, i.e. on a un
hom naturel $A\unsure{-|\to} A\otimes A$, comme
on vérifie aussitôt (on est ramené
au cas de $\Omega=L_{k}(V)$, $V$ de dim
finie sur $k$).
\inserted{de plus \unsure{$A$ est} réunion \unsure{filtrante} de ses sous-coalgèbres de dim finie}
Alors, les $\Omega$-modules
unitaires \deleted{\unsure{de dim finie}} \deleted{\unsure{sont exactement les}} de
